\section{格把偏序变成代数运算}
\label{sec:04-Discrete-Mathematics/08-Rings-and-Lattices/02}

编译器读到一个条件表达式，两个分支分别返回 \texttt{Int32} 与 \texttt{Float64}，它必须当场给整个表达式定一个类型。这个类型要同时容得下两支的结果，所以它得是两者共同的上界。可是共同上界通常有一堆——\texttt{Object} 也是，任何足够宽的类型都是。全都给 \texttt{Object}，类型检查就废了，后面那句加法再也查不出错。编译器要的是所有共同上界里最小的那个。

\hyperref[sec:01-Mathematical-Language/02-Sets-Relations-Functions/03]{``偏序与可比较性''}给的关系只回答一件事：\texttt{Int32} 是不是 \texttt{Float64} 的子类型。它不承诺``最小的共同上界''存在，更不承诺唯一。一旦这个承诺补上，事情的性质就变了：``取共同上界的最小者''不再是每次都要现搜一遍的查询，而是一个二元运算，可以像加法乘法那样连着写、结合着算。

\subsection{把上确界和下确界当成运算}

偏序集 \((L,\le)\) 里，若任意两个元素 \(x,y\) 都有最大的共同下界 \(x\wedge y\)（读作 meet）和最小的共同上界 \(x\vee y\)（读作 join），就称 \(L\) 是格。两个运算的名字听着抽象，具体化之后全是老熟人：幂集按包含排序时 \(A\wedge B=A\cap B\) 而 \(A\vee B=A\cup B\)；正整数按整除排序时 \(a\wedge b=\gcd(a,b)\) 而 \(a\vee b=\operatorname{lcm}(a,b)\)。集合的交并与数论里的 \(\gcd\)、\(\operatorname{lcm}\)，从这个角度看是同一对运算穿了不同的衣服。

\begin{center}
\begin{tikzpicture}[x=1cm,y=1cm,font=\small]
  \draw[black!55] (0,0) -- (-1.2,1.5);
  \draw[black!55] (0,0) -- (1.2,1.5);
  \draw[black!55] (-1.2,1.5) -- (-1.2,3.0);
  \draw[black!55] (-1.2,1.5) -- (1.2,3.0);
  \draw[black!55] (1.2,1.5) -- (1.2,3.0);
  \draw[black!55] (-1.2,3.0) -- (0,4.5);
  \draw[black!55] (1.2,3.0) -- (0,4.5);
  \fill[white] (0,0) circle (0.28);
  \draw[black!60] (0,0) circle (0.28);
  \node at (0,0) {1};
  \fill[blue!25] (-1.2,1.5) circle (0.28);
  \draw[black!60] (-1.2,1.5) circle (0.28);
  \node at (-1.2,1.5) {2};
  \fill[white] (1.2,1.5) circle (0.28);
  \draw[black!60] (1.2,1.5) circle (0.28);
  \node at (1.2,1.5) {3};
  \fill[orange!35] (-1.2,3.0) circle (0.28);
  \draw[black!60] (-1.2,3.0) circle (0.28);
  \node at (-1.2,3.0) {4};
  \fill[orange!35] (1.2,3.0) circle (0.28);
  \draw[black!60] (1.2,3.0) circle (0.28);
  \node at (1.2,3.0) {6};
  \fill[blue!25] (0,4.5) circle (0.28);
  \draw[black!60] (0,4.5) circle (0.28);
  \node at (0,4.5) {12};
  \node[right,black!75,font=\footnotesize] at (0.45,4.5) {\(4\vee6=\operatorname{lcm}(4,6)=12\)};
  \node[left,black!75,font=\footnotesize] at (-1.75,1.5) {\(4\wedge6=\gcd(4,6)=2\)};
  \node[right,black!60,font=\footnotesize] at (1.65,3.0) {两个元素};
  \node[below,black!60,font=\footnotesize] at (0,-0.45) {整除关系向上递增};
\end{tikzpicture}

\emph{12 的因子按整除关系画成 Hasse 图，连线表示``下面那个整除上面那个''。橙色两点的共同上界有 12 一个，共同下界有 1 和 2 两个，取最大者得 2。注意 4 与 6 本身不可比——格不要求任意两元素能比较，只要求它们的上下确界存在。}
\end{center}

图上的操作可以照着做一遍：从两个点各自向上追，第一次汇合的地方是 join；向下追，第一次汇合的地方是 meet。这个``第一次汇合''的措辞不严谨但直觉对，正式的说法仍然是共同上界中的最小者。

\subsection{Hasse 图画得出来，不代表它是格}

回到开头那个类型问题。类 \(A\) 和类 \(B\) 各自实现了接口 \(I\) 与 \(J\)，而 \(I\) 与 \(J\) 之间毫无关系。那么 \(A\vee B\) 是谁？\(I\) 是共同上界，\(J\) 也是，两者互不比较，没有哪一个更小。最小上界不存在。

\begin{center}
\begin{tikzpicture}[x=1cm,y=1cm,font=\small]
  \draw[black!55] (-1.3,0) -- (-1.3,2.2);
  \draw[black!55] (-1.3,0) -- (1.3,2.2);
  \draw[black!55] (1.3,0) -- (-1.3,2.2);
  \draw[black!55] (1.3,0) -- (1.3,2.2);
  \fill[white] (-1.3,0) circle (0.3);
  \draw[black!60] (-1.3,0) circle (0.3);
  \node at (-1.3,0) {\(A\)};
  \fill[white] (1.3,0) circle (0.3);
  \draw[black!60] (1.3,0) circle (0.3);
  \node at (1.3,0) {\(B\)};
  \fill[red!22] (-1.3,2.2) circle (0.3);
  \draw[black!60] (-1.3,2.2) circle (0.3);
  \node at (-1.3,2.2) {\(I\)};
  \fill[red!22] (1.3,2.2) circle (0.3);
  \draw[black!60] (1.3,2.2) circle (0.3);
  \node at (1.3,2.2) {\(J\)};
  \node[right,black!75,font=\footnotesize] at (1.85,2.2) {两个极小上界，互不可比};
  \node[right,black!60,font=\footnotesize] at (1.85,0) {\(A\vee B\) 无定义};
  \fill[black!8] (0,3.9) circle (0.3);
  \draw[black!45,dashed] (0,3.9) circle (0.3);
  \node[black!60] at (0,3.9) {\(\top\)};
  \draw[black!45,dashed] (-1.3,2.5) -- (0,3.6);
  \draw[black!45,dashed] (1.3,2.5) -- (0,3.6);
  \node[right,black!60,font=\footnotesize] at (0.45,3.9) {补一个顶元素才勉强补回 join};
\end{tikzpicture}

\emph{一个不是格的偏序。\(A\) 与 \(B\) 有两个极小的共同上界，谁也不比谁小，join 无从取值。真实的类型系统要么像 Java 那样补一个 \texttt{Object} 兜底，要么引入交类型让 \(I\) 与 \(J\) 的组合本身成为一个类型；两条路都是在动手把这个偏序补成格。}
\end{center}

所以画完 Hasse 图别急着说``这是个格''。要检查的是存在性与唯一性：共同上界这一堆里，有没有一个被其余全部包住？多个极小上界并列，或者干脆没有上界，格的资格就没了。上一节的环也有类似的体检项——那里问的是有没有零因子，这里问的是上下确界在不在。

\subsection{四条恒等式足以把序找回来}

meet 与 join 各自交换、结合、幂等（\(x\wedge x=x\)），彼此之间还满足吸收律 \(x\wedge(x\vee y)=x\) 与 \(x\vee(x\wedge y)=x\)。这些都能从上下确界的定义直接读出，验证过程没什么可讲的。有意思的是反过来：只要一对运算满足这几条恒等式，令

\[
x\le y
\quad\Longleftrightarrow\quad
x\wedge y=x,
\]

就能把偏序造回来，而且造出来的序恰好使这对运算成为它的 meet 和 join。序和运算是同一件事的两种写法，你可以从任一头进入。

这个双向通道在分布式系统里被用得很直白。允许多副本各自写入、之后再合并的数据结构（CRDT），核心要求就是合并函数满足交换、结合、幂等三条。满足了，就意味着状态空间上暗含一个偏序，而合并等于取 join：无论消息以什么顺序到达、重复多少次，各副本都会爬向同一个上确界。收敛性不是靠协调协议换来的，是靠代数律换来的。

反例同样值得记住。真实系统里的权限合并常常带例外规则——某个组的成员即便属于另一个组也要被降权。这类上下文依赖一进来，合并往往不再结合，先合 \(a,b\) 再合 \(c\) 与先合 \(b,c\) 再合 \(a\) 结果不同，那就不是格，前面那套保证一条都不剩。

\subsection{分配律与补元不是白送的}

格给了两个运算，但没有承诺它们像集合的交并那样互相分配。若

\[
x\wedge(y\vee z)=(x\wedge y)\vee(x\wedge z)
\]

及其对偶式成立，称为分配格；在此之上再要求有最小元 0、最大元 1，且每个 \(x\) 配一个补元 \(x'\) 满足 \(x\wedge x'=0\) 与 \(x\vee x'=1\)，得到的就是布尔代数。\hyperref[sec:04-Discrete-Mathematics/07-Boolean-Algebra-and-Automata/01]{``布尔代数与逻辑化简''}里的 AND、OR、NOT，在这个视角下正是分配格上的 meet、join 与补。

漏掉哪一条，逻辑就会跟着变形。补元在一般格里既不保证存在也不保证唯一；分配律更是随时会塌——线性子空间按包含排序构成格，meet 是交、join 是线性和，它不分配。下一节会拿这个例子说明推理规则怎样依赖底层格的形状。

\subsection{完全格装得下不动点}

前面只要求两个元素有 meet 和 join。把``两个''换成``任意子集''，得到完全格：任何一堆元素，哪怕无穷多，都有上确界和下确界。幂集是完全格，有限格自动是完全格；有理数按大小排序则不是，一个有上界的有理数集合，上确界可能是无理数，跑到集合外面去了。这个升级看着技术，回报却很大。

\begin{theorem}[Knaster--Tarski]
完全格 \(L\) 上的单调映射 \(T\)（\(x\le y\) 蕴含 \(T(x)\le T(y)\)）必有不动点，且全体不动点在原序下仍构成完全格。特别地，最小不动点与最大不动点都存在。
\end{theorem}

证明是构造性的一行——取所有满足 \(T(x)\le x\) 的元素之下确界——但真正该记住的是它保证了什么。递归定义、程序数据流分析、Datalog 的求值、类型推断的迭代，做的都是同一件事：从最小的信息开始，反复施加一条单调的更新规则，直到不再变化。这类过程的``答案''被定义为最小不动点，而定理说的正是这个答案一定存在，不需要每换一个规则系统就重新论证一次。注意 \(T\) 只要求单调，不要求连续或收缩；条件之弱是它好用的原因。

这里的``收敛''与度量空间里的收敛不是一回事。那边说的是距离逐步缩小，这边说的是沿偏序单调爬升、信息只增不减。两者都会停在某处，理由完全不同。存在性归存在性，迭代要在有限步内摸到那个不动点还需要额外条件，格得足够矮才行。

下一节把这三样东西——单调映射、最小不动点、格的高度——放进一个具体场合：程序分析要在不运行程序的前提下给出可靠结论，靠的是两个格之间一对方向相反的映射。
